\section{Results}
\label{sec:results}

\begin{table}[t]
\centering
\small
\caption{Steady-state performance at 80\% offered load, averaged over five runs.}
\label{tab:results}
\begin{tabular}{lrrrr}
\toprule
\textbf{Policy} & \textbf{p50} & \textbf{p99} & \textbf{Tput.} & \textbf{Ovh.} \\
 & \textbf{(ms)} & \textbf{(ms)} & \textbf{(kreq/s)} & \textbf{(\%)} \\
\midrule
Least-Loaded    & 6.1  & 88.4 & 41.2 & 0.4 \\
PowerOfTwo      & 5.4  & 61.7 & 46.8 & 0.6 \\
Tempo (ours)    & \textbf{4.8} & \textbf{47.5} & \textbf{52.1} & 1.1 \\
\bottomrule
\end{tabular}
\end{table}

Table~\ref{tab:results} summarizes steady-state behavior at 80\% offered
load, the region where the three schedulers diverge most. Tempo reduces
p99 latency by 41\% relative to Least-Loaded and by 23\% relative to
PowerOfTwo, while also sustaining higher throughput, because its packing
decisions free up headroom that the other two policies waste on
already-quiet nodes. Tempo's scheduler-side CPU overhead is roughly double
that of PowerOfTwo, almost entirely attributable to the recursive
least-squares update in Equation~\ref{eq:predictor}; at 1.1\% of one core
per scheduler instance, we consider this cost negligible relative to the
latency win.

Figure~\ref{fig:latency-load} shows p99 latency as offered load increases
from 30\% to 90\% of fleet capacity. Below 50\% load the three policies are
statistically indistinguishable, since no policy is forced to make a hard
trade-off yet. Above roughly 60\% load, Least-Loaded's latency grows
sharply as its conservative placement runs out of genuinely idle nodes,
while Tempo continues to track a much flatter curve by anticipating
contention before utilization-based signals would flag it.

\begin{figure}[t]
\centering
\begin{tikzpicture}
\begin{axis}[
    width=0.98\linewidth,
    height=4.6cm,
    xlabel={Offered load (\% of fleet capacity)},
    ylabel={p99 latency (ms)},
    xmin=30, xmax=90,
    ymin=0, ymax=140,
    legend pos=north west,
    legend style={font=\footnotesize, draw=gray!40},
    label style={font=\footnotesize},
    tick label style={font=\footnotesize},
    grid=major,
    grid style={gray!20},
    axis line style={gray!70},
]
\addplot[color=usenixblue!60,mark=square*,mark size=1.4pt,thick]
  coordinates {(30,9)(40,11)(50,15)(60,24)(70,45)(80,88)(90,131)};
\addlegendentry{Least-Loaded}

\addplot[color=accentgold,mark=triangle*,mark size=1.6pt,thick]
  coordinates {(30,8)(40,10)(50,13)(60,19)(70,31)(80,62)(90,104)};
\addlegendentry{PowerOfTwo}

\addplot[color=usenixblue,mark=*,mark size=1.4pt,thick]
  coordinates {(30,8)(40,9)(50,11)(60,15)(70,22)(80,48)(90,69)};
\addlegendentry{Tempo (ours)}
\end{axis}
\end{tikzpicture}
\caption{Tail latency versus offered load. Tempo's predictive placement
delays the onset of latency inflation relative to both baselines.}
\label{fig:latency-load}
\end{figure}

We also measured sensitivity to the adaptation rate of
$\alpha(t),\beta(t),\gamma(t)$. Disabling adaptation (fixing the weights at
their converged values from a single workload mix) degrades p99 latency by
roughly 12\% when the workload composition is changed mid-run, confirming
that the online adjustment described in Section~\ref{sec:method}, not just
the predictor itself, contributes to Tempo's robustness across workloads.
